\numoptsfalse
\yearpaper{2019}{Version A-1 \textperiodcentered{} 60 Questions \textperiodcentered{} 60 Marks \textperiodcentered{} 80 Minutes}{}
\begin{multicols}{2}

\QQ{In a cyclotron a charged particle\ntag{X}{}{}{}{}{cyclotron}}
{\opt{undergoes acceleration all the time}{speeds up between the dees because of the magnetic field.}{speeds up in dee}{slows down within a dee and speeds up between dees}}

\QQ{The number of turns in a coil of Galvanometer is tripled, then\ntag{D}{131}{leph104}{25}{}{Class XII Part 1 \textperiodcentered{} Ch 4}}
{\opt{Voltage sensitivity increases 3 times and current sensitivity remains constant}{Voltage sensitivity remains constant and current sensitivity increases 3 times}{Both voltage and current sensitivity remains constant}{Both voltage and current sensitivity decreases by 33\%}}

\QQ{A circular current loop of magnetic moment M is in an arbitrary orientation in an external uniform magnetic field $\vec B$. The work done to rotate the loop by $30^\circ$ about an axis perpendicular to its plane is\ntag{A}{126}{leph104}{20}{ \textperiodcentered{} Exemplar 4.6}{Class XII Part 1 \textperiodcentered{} Ch 4}\nsrc{None found in NEET/AIPMT/AIEEE/JEE papers, 2005--2026 (NCERT concept application)}}
{\opt{MB}{$\sqrt3\,\dfrac{MB}{2}$}{$\dfrac{MB}{2}$}{Zero}}

\QQ{In a permanent magnet at room temperature\ntag{D}{149}{leph105}{14}{ \textperiodcentered{} Exemplar 5.3}{Class XII Part 1 \textperiodcentered{} Ch 5}}
{\opt{magnetic moment of each molecule is zero.}{the individual molecules have non zero magnetic moment which are all perfectly aligned.}{domains are partially aligned.}{domains are all perfectly aligned.}}

\QQ{Coersivity of a magnet where the ferromagnet gets completely demagnetized is $3\times10^3\unit{A\,m^{-1}}$. The minimum current required to be passed in a solenoid having 1000 turns per metre, so that the magnet gets completely demagnetized when placed inside the solenoid is\ntag{X}{}{}{}{}{coercivity (hysteresis)}}
{\opt{$30\unit{mA}$}{$60\unit{mA}$}{$3\unit{A}$}{$6\unit{A}$}}

\QQ{Which one of the following nuclei has shorter mean life?
\fig{\begin{tikzpicture}[scale=0.8]
\draw[->] (0,0) -- (2.8,0) node[right,lbl]{$t$}; \draw[->] (0,0) -- (0,2.0) node[above,lbl]{$\dfrac{dN}{dt}$};
\draw (0,1.4) .. controls (0.8,1.0) and (1.6,0.8) .. (2.4,0.7) node[right,lbl]{C};
\draw (0,1.4) .. controls (0.6,0.8) and (1.4,0.5) .. (2.2,0.42) node[right,lbl]{B};
\draw (0,1.4) .. controls (0.5,0.6) and (1.2,0.3) .. (2.0,0.22) node[right,lbl]{A};
\end{tikzpicture}}\ntag{X}{}{}{}{}{radioactive decay law (mean life)}}
{\opt{A}{B}{C}{Same for all}}

\QQ{The conductivity of semiconductor increases with increase in temperature because\ntag{D}{91}{leph103}{11}{ \textperiodcentered{} Exemplar 14.1}{Class XII Part 1 \textperiodcentered{} Ch 3}}
{\opt{number density of charge carriers increases}{relaxation time increases}{both number density of charge carriers and relaxation time increase}{number density of current carriers increases, relaxation time decreases but effect of decrease in relaxation time is much less than increase in number density}}

\QQ{For a transistor amplifier, the voltage gain\ntag{X}{}{}{}{}{transistor amplifier}}
{\opt{remains constant for all frequencies}{is high at high and low frequencies and constant in the middle frequency range}{is low at high and low frequencies and constant at mid frequencies}{constant at high frequencies and low at low frequencies}}

\QQ{In the following circuit, what are P and Q?
\fig{\begin{circuitikz}[scale=0.7,transform shape]
\draw (0,1.4) node[nor port](G1){}; \draw (0,-0.4) node[nor port](G2){};
\draw (G1.in 1) -- ++(-0.6,0) node[ocirc]{} node[left]{1};
\draw (G2.in 2) -- ++(-0.6,0) node[ocirc]{} node[left]{0};
\draw (G1.out) -- ++(0.6,0) node[ocirc]{} node[right]{P}; \draw (G2.out) -- ++(0.6,0) node[ocirc]{} node[right]{Q};
\draw (G1.out) ++(0.2,0) -- ++(0,-0.4) -- ($(G2.in 1)+(-0.3,0.5)$) -- ($(G2.in 1)+(-0.3,0)$) -- (G2.in 1);
\draw (G2.out) ++(0.2,0) -- ++(0,0.4) -- ($(G1.in 2)+(-0.3,-0.5)$) -- ($(G1.in 2)+(-0.3,0)$) -- (G1.in 2);
\end{circuitikz}}\ntag{X}{}{}{}{}{digital electronics (logic gates)}}
{\opt{$P = 0$, $Q = 0$}{$P = 1$, $Q = 0$}{$P = 0$, $Q = 1$}{$P = 1$, $Q = 1$}}

\QQ{An antenna uses electromagnetic waves of frequency $5\unit{MHz}$. For proper working, the size of the antenna should be\ntag{X}{}{}{}{}{antenna size (communication systems)}}
{\opt{$15\unit{m}$}{$300\unit{m}$}{$15\unit{km}$}{$3\unit{km}$}}

\QQ{If P, Q and R are physical quantities having different dimensions, which of the following combinations can never be a meaningful quantity?\ntag{D}{8}{keph101}{8}{}{Class XI Part 1 \textperiodcentered{} Ch 1}}
{\opt{$\dfrac{P-Q}{R}$}{$PQ-R$}{$\dfrac{PQ}{R}$}{$\dfrac{PR-Q^2}{R}$}}

\QQ{The given graph shows the variation of velocity ($v$) with position ($x$) for a particle moving along a straight line
\fig{\gaxes{$x$}{$v$}{\draw (0,1.5) -- (1.6,0); \node[lbl,below left] at (0,0) {O};}}
Which of the following graph shows the variation of acceleration ($a$) with position ($x$)?\ntag{A}{16}{keph102}{4}{}{Class XI Part 1 \textperiodcentered{} Ch 2}\nsrc{None found in NEET/AIPMT/AIEEE/JEE papers, 2005--2026 (NCERT concept application)}}
{\optfig{\gaxesm{$x$}{$a$}{\draw (0,0.9) -- (1.6,0);}}{\gaxesm{$x$}{$a$}{\draw (0,-0.6) -- (1.9,-0.6);}}{\gaxesm{$x$}{$a$}{\draw (0,-0.9) -- (1.6,0);}}{\gaxesm{$x$}{$a$}{\draw (0,0) -- (1.6,0.9);}}}

\QQ{The trajectory of a projectile projected from origin is given by the equation $y = x-\dfrac{2x^2}{5}$. The initial velocity of the projectile is\ntag{A}{39}{keph103}{13}{}{Class XI Part 1 \textperiodcentered{} Ch 3}\nsrc{None found in NEET/AIPMT/AIEEE/JEE papers, 2005--2026 (NCERT concept application)}}
{\opt{$\dfrac25\unit{m\,s^{-1}}$}{$5\unit{m\,s^{-1}}$}{$25\unit{m\,s^{-1}}$}{$\dfrac52\unit{m\,s^{-1}}$}}

\QQ{An object with mass $5\unit{kg}$ is acted upon by a force, $\vec F = (-3\hat i+4\hat j)\unit{N}$. If its initial velocity at $t = 0$ is $\vec v = (6\hat i-12\hat j)\unit{m\,s^{-1}}$, the time at which it will just have a velocity along $y$-axis is\ntag{A}{53}{keph104}{5}{}{Class XI Part 1 \textperiodcentered{} Ch 4}\nsrc{None found in NEET/AIPMT/AIEEE/JEE papers, 2005--2026 (NCERT concept application)}}
{\opt{$5\unit{s}$}{$10\unit{s}$}{$2\unit{s}$}{$15\unit{s}$}}

\QQ{During inelastic collision between two objects, which of the following quantity always remains conserved?\ntag{D}{84}{keph105}{14}{}{Class XI Part 1 \textperiodcentered{} Ch 5}}
{\opt{Total kinetic energy}{Total mechanical energy}{Total linear momentum}{Speed of each body}}

\QQ{In Rutherford experiment, for head-on collision of $\alpha$-particles with a gold nucleus, the impact parameter is\ntag{D}{294}{leph204}{5}{}{Class XII Part 2 \textperiodcentered{} Ch 12}}
{\opt{zero}{of the order of $10^{-14}\unit{m}$}{of the order of $10^{-10}\unit{m}$}{of the order of $10^{-6}\unit{m}$}}

\QQ{Frequency of revolution of an electron revolving in $n^{\text{th}}$ orbit of H-atom is proportional to\ntag{A}{299}{leph204}{10}{}{Class XII Part 2 \textperiodcentered{} Ch 12}\nsrc{None found in NEET/AIPMT/AIEEE/JEE papers, 2005--2026 (NCERT concept application)}}
{\opt{$\dfrac1{n^2}$}{$n$}{$n$ independent of $n$}{$\dfrac1{n^3}$}}

\QQ{A hydrogen atom in ground state absorbs $10.2\unit{eV}$ of energy. The orbital angular momentum of the electron is increased by\ntag{A}{300}{leph204}{11}{}{Class XII Part 2 \textperiodcentered{} Ch 12}\nsrc{None found in NEET/AIPMT/AIEEE/JEE papers, 2005--2026 (NCERT concept application)}}
{\opt{$1.05\times10^{-34}\unit{J\,s}$}{$2.11\times10^{-34}\unit{J\,s}$}{$3.16\times10^{-34}\unit{J\,s}$}{$4.22\times10^{-34}\unit{J\,s}$}}

\QQ{The end product of decay of ${}_{90}\mathrm{Th}^{232}$ is ${}_{82}\mathrm{Pb}^{208}$. The number of $\alpha$ and $\beta$ particles emitted are respectively\ntag{X}{}{}{}{}{alpha- and beta-decay}}
{\opt{3, 3}{6, 4}{6, 0}{4, 6}}

\QQ{Two protons are kept at a separation of $10\unit{nm}$. Let $F_n$ and $F_e$ be the nuclear force and the electromagnetic force between them\ntag{D}{313}{leph205}{8}{}{Class XII Part 2 \textperiodcentered{} Ch 13}}
{\opt{$F_e = F_n$}{$F_e\gg F_n$}{$F_e\ll F_n$}{$F_e$ and $F_n$ differ only slightly}}

\QQ{A transparent medium shows relation between $i$ and $r$ as shown. If the speed of light in vacuum is $c$ the Brewster angle for the medium is
\fig{\gaxes{$\sin i$}{$\sin r$}{\draw (0,0) -- (1.9,1.1); \draw (0.6,0) arc (0:30:0.6); \node[lbl] at (0.85,0.2) {$30^\circ$};}}\ntag{X}{}{}{}{}{Brewster's law}}
{\opt{$30^\circ$}{$45^\circ$}{$60^\circ$}{$90^\circ$}}

\QQ{In Young's double slit experiment, using monochromatic light of wavelength $\lambda$, the intensity of light at a point on the screen where path difference is $\lambda$ is K units. The intensity of light at a point where path difference is $\dfrac\lambda3$ is\ntag{D}{273}{leph202}{19}{ \textperiodcentered{} Exercise 10.5}{Class XII Part 2 \textperiodcentered{} Ch 10}}
{\opt{K}{$\dfrac K4$}{4K}{2K}}

\QQ{Due to Doppler's effect the shift in wavelength observed is $0.1\unit{\text{\AA}}$ for a star producing wavelength $6000\unit{\text{\AA}}$. Velocity of recession of the star will be\ntag{X}{}{}{}{}{Doppler effect for light}}
{\opt{$25\unit{km/s}$}{$10\unit{km/s}$}{$5\unit{km/s}$}{$20\unit{km/s}$}}

\QQ{An electron is moving with an initial velocity $\vec V = V_0\hat i$ and is in a uniform magnetic field $\vec B = B_0\hat j$. Then its de Broglie wavelength\ntag{A}{285}{leph203}{12}{}{Class XII Part 2 \textperiodcentered{} Ch 11}\nsrc{None found in NEET/AIPMT/AIEEE/JEE papers, 2005--2026 (NCERT concept application)}}
{\opt{remains constant}{increases with time}{decreases with time}{increase and decreases periodically}}

\QQ{Light of certain frequency and intensity incident on a photosensitive material causes photoelectric effect. If both the frequency and intensity are doubled, the photoelectric saturation current becomes\ntag{D}{279}{leph203}{6}{}{Class XII Part 2 \textperiodcentered{} Ch 11}}
{\opt{quadrupled}{doubled}{halved}{unchanged}}

\QQ{A magnetic needle has a magnetic moment of $5\times10^{-2}\unit{A\,m^2}$ and moment of inertia $8\times10^{-6}\unit{kg\,m^2}$. It has a period of oscillation of $2\unit{s}$ in a magnetic field $\vec B$. The magnitude of magnetic field is approximately\ntag{X}{}{}{}{}{oscillation period of a magnetic needle}}
{\opt{$1.6\times10^{-4}\unit{T}$}{$0.4\times10^{-4}\unit{T}$}{$3.2\times10^{-4}\unit{T}$}{$0.8\times10^{-4}\unit{T}$}}

\QQ{A toroid has 500 turns per metre length. If it carries a current of $2\unit{A}$, the magnetic energy density inside the toroid is\ntag{X}{}{}{}{}{toroid}}
{\opt{$0.628\unit{J/m^3}$}{$0.314\unit{J/m^3}$}{$6.28\unit{J/m^3}$}{$3.14\unit{J/m^3}$}}

\QQ{Consider the situation given in figure. The wire AB is slid on the fixed rails with a constant velocity. If the wire AB is replaced by a semicircular wire, the magnitude of the induced current will
\fig{\begin{tikzpicture}[scale=0.55]
\foreach \x in {0,1,...,6} \foreach \y in {0,1,...,5} \node[font=\scriptsize] at (\x,\y) {$\times$};
\draw (2.5,4.4) -- (6.6,4.4) -- (6.6,0.6) -- (2.5,0.6);
\draw[line width=1pt] (2.5,4.6) -- (2.5,0.4); \node[lbl] at (2.85,4.75) {A}; \node[lbl] at (2.85,0.25) {B};
\draw[->] (2.3,3.0) -- (1.2,3.0) node[above,lbl]{$v$};
\draw[decorate,decoration={zigzag,segment length=3pt,amplitude=2pt}] (6.6,3.2) -- (6.6,1.8);
\end{tikzpicture}}\ntag{A}{163}{leph106}{10}{}{Class XII Part 1 \textperiodcentered{} Ch 6}\nsrc{None found in NEET/AIPMT/AIEEE/JEE papers, 2005--2026 (NCERT concept application)}}
{\opt{increase}{remain same}{decrease}{increase or decrease depending on whether the semicircle bulges towards the resistance or away from it}}

\QQ{The frequency of an alternating current is $50\unit{Hz}$. What is the minimum time taken by current to reach its peak value from rms value?\ntag{A}{180}{leph107}{4}{}{Class XII Part 1 \textperiodcentered{} Ch 7}\nsrc{None found in NEET/AIPMT/AIEEE/JEE papers, 2005--2026 (NCERT concept application)}}
{\opt{$5\times10^{-3}\unit{s}$}{$2.5\times10^{-3}\unit{s}$}{$0.02\unit{s}$}{$10\times10^{-3}\unit{s}$}}

\QQ{The readings of ammeter and voltmeter in the following circuit are respectively
\fig{\begin{circuitikz}[scale=0.6,transform shape,font=\Large]
\draw (0,2.4) to[L,l=L,a={$V_L=50\,\mathrm V$}] (2.6,2.4) to[C,l=C,a={$V_C=50\,\mathrm V$}] (4.8,2.4) to[R,l={$R=100\,\Omega$}] (7.2,2.4) -- (7.2,0) to[sV,l={220 V, 50 Hz}] (0,0) to[rmeter,t=A] (0,2.4);
\draw (4.8,2.4) -- (4.8,1.4) to[rmeter,t=V] (7.2,1.4) -- (7.2,2.4);
\end{circuitikz}}\ntag{A}{188}{leph107}{12}{}{Class XII Part 1 \textperiodcentered{} Ch 7}\nsrc{None found in NEET/AIPMT/AIEEE/JEE papers, 2005--2026 (NCERT concept application)}}
{\opt{1.2\,A, 120\,V}{1.5\,A, 100\,V}{2.7\,A, 220\,V}{2.2\,A, 220\,V}}

\QQ{Two particles which are initially at rest move towards each other under the action of their mutual attraction. If their speeds are $v$ and $2v$ at any instant, then the speed of center of mass of the system is,\ntag{D}{99}{keph106}{8}{}{Class XI Part 1 \textperiodcentered{} Ch 6}}
{\opt{$2v$}{Zero}{$1.5v$}{$v$}}

\QQ{A particle is moving uniformly along a straight line as shown in the figure. During the motion of the particle from A to B, the angular momentum of the particle about `O'
\fig{\begin{tikzpicture}[scale=0.8]
\draw[->] (-1.0,0) -- (2.2,0) node[right,lbl]{$x$}; \draw[->] (0,-0.8) -- (0,1.7) node[above,lbl]{$y$};
\node[lbl,below left] at (0,0) {O};
\draw (-0.9,-0.4) -- (1.0,1.5); \draw[->] (0.4,0.9) -- (0.6,1.1);
\fill (-0.3,0.2) circle (1.5pt) node[left,lbl]{A}; \fill (0.2,0.7) circle (1.5pt) node[below right,lbl]{P}; \node[lbl,left] at (0.75,1.3) {B};
\end{tikzpicture}}\ntag{A}{109}{keph106}{18}{}{Class XI Part 1 \textperiodcentered{} Ch 6}\nsrc{None found in NEET/AIPMT/AIEEE/JEE papers, 2005--2026 (NCERT concept application)}}
{\opt{increases}{decreases}{remains constant}{first increases then decreases}}

\QQ{A satellite is orbiting close to the earth and has a kinetic energy K. The minimum extra kinetic energy required by it to just overcome the gravitation pull of the earth is\ntag{A}{135}{keph107}{9}{}{Class XI Part 1 \textperiodcentered{} Ch 7}\nsrc{None found in NEET/AIPMT/AIEEE/JEE papers, 2005--2026 (NCERT concept application)}}
{\opt{K}{2K}{$\sqrt3$K}{$2\sqrt2$K}}

\QQ{A wire is stretched such that its volume remains constant. The Poission's ratio of the material of the wire is\ntag{A}{174}{keph201}{8}{}{Class XI Part 2 \textperiodcentered{} Ch 8}\nsrc{None found in NEET/AIPMT/AIEEE/JEE papers, 2005--2026 (NCERT concept application)}}
{\opt{0.50}{$-0.50$}{0.25}{$-0.25$}}

\QQ{A cylindrical container containing water has a small hole at height of $H = 8\unit{cm}$ from the bottom and at a depth of $2\unit{cm}$ from the top surface of the liquid. The maximum horizontal distance travelled by the water before it hits the ground ($x$) is
\fig{\begin{tikzpicture}[scale=0.7]
\draw (-0.4,0) -- (3.6,0); \draw (0,0) -- (0,2.6) (2.6,0) -- (2.6,2.6);
\fill[gray!45] (0.03,0.03) rectangle (2.57,2.4);
\draw[densely dashed] (0,1.9) -- (2.6,1.9); \draw[->] (2.6,1.9) -- (2.95,1.9);
\draw[<->] (-0.2,1.9) -- (-0.2,2.4) node[midway,left,lbl]{$h$}; \draw[<->] (-0.2,0) -- (-0.2,1.9) node[midway,left,lbl]{$H$};
\end{tikzpicture}}\ntag{A}{189}{keph202}{10}{}{Class XI Part 2 \textperiodcentered{} Ch 9}\nsrc{None found in NEET/AIPMT/AIEEE/JEE papers, 2005--2026 (NCERT concept application)}}
{\opt{$8\unit{cm}$}{$4\sqrt2\unit{cm}$}{$4\unit{cm}$}{$6\unit{cm}$}}

\QQ{An inductor of inductance L and resistor R are joined together in series and connected by a source of frequency $\omega$. The power dissipated in the circuit is\ntag{A}{191}{leph107}{15}{}{Class XII Part 1 \textperiodcentered{} Ch 7}\nsrc{None found in NEET/AIPMT/AIEEE/JEE papers, 2005--2026 (NCERT concept application)}}
{\opt{$\dfrac{R^2+\omega^2L^2}{V}$}{$\dfrac{V^2R}{R^2+\omega^2L^2}$}{$\dfrac{V}{R^2+\omega^2L^2}$}{$\dfrac{V^2R}{\sqrt{R^2+\omega^2L^2}}$}}

\QQ{An electromagnetic wave is travelling in $x$-direction with electric field vector given by, $\vec E_y = E_0\sin(kx-\omega t)\,\hat j$. The correct expression for magnetic field vector is\ntag{D}{206}{leph108}{6}{}{Class XII Part 1 \textperiodcentered{} Ch 8}}
{\opt{$\vec B_y = E_0C\sin(kx-\omega t)\,\hat j$}{$\vec B_z = E_0C\sin(kx-\omega t)\,\hat k$}{$\vec B_y = \dfrac{E_0}{C}\sin(kx-\omega t)\,\hat j$}{$\vec B_z = \dfrac{E_0}{C}\sin(kx-\omega t)\,\hat k$}}

\QQ{The phenomenon involved in the reflection of radio-waves by ionosphere is similar to\ntag{A}{246}{leph201}{26}{}{Class XII Part 2 \textperiodcentered{} Ch 9}\nsrc{None found in NEET/AIPMT/AIEEE/JEE papers, 2005--2026 (NCERT concept application)}}
{\opt{reflection of light by plane mirror}{total internal reflection of light in air during a mirage}{dispersion of light by water molecules during the formation of a rainbow}{scattering of light by air particles}}

\QQ{A point object is moving uniformly towards the pole of a concave mirror of focal length $25\unit{cm}$ along its axis as shown below. The speed of the object is $1\unit{m\,s^{-1}}$. At $t = 0$, the distance of the object from the mirror is $50\unit{cm}$. The average velocity of the image formed by the mirror between time $t = 0$ and $t = 0.25\unit{s}$ is:
\fig{\begin{tikzpicture}[scale=0.8]
\draw (-0.4,0) -- (3.4,0); \node[lbl,above] at (-0.4,0) {$t=0$};
\fill (0.3,0) circle (1.4pt) node[below,lbl]{O}; \draw[->] (0.3,0.25) -- (1.0,0.25) node[right,lbl]{$V = 1\unit{m\,s^{-1}}$};
\draw[line width=1pt] (2.75,-0.8) arc (-25:25:1.9); \foreach \t in {-22,-16,...,22} \draw ($(2.75,-0.8)+(-25:0)$) ;
\node[lbl,right] at (2.95,0.15) {P};
\draw[<->] (0.3,-0.9) -- (2.95,-0.9) node[midway,below,lbl]{50 cm};
\end{tikzpicture}}\ntag{A}{224}{leph201}{4}{}{Class XII Part 2 \textperiodcentered{} Ch 9}\nsrc{None found in NEET/AIPMT/AIEEE/JEE papers, 2005--2026 (NCERT concept application)}}
{\opt{$40\unit{cm\,s^{-1}}$}{$20\unit{cm\,s^{-1}}$}{Zero}{Infinity}}

\QQ{A certain prism is found to produce a minimum deviation of $38^\circ$. It produces a deviation of $44^\circ$ when the angle of incidence is either $42^\circ$ or $62^\circ$. What is the angle of incidence when it is undergoing minimum deviation?\ntag{A}{240}{leph201}{20}{}{Class XII Part 2 \textperiodcentered{} Ch 9}\nsrc{None found in NEET/AIPMT/AIEEE/JEE papers, 2005--2026 (NCERT concept application)}}
{\opt{$30^\circ$}{$40^\circ$}{$49^\circ$}{$60^\circ$}}

\QQ{An aluminium sphere is dipped into water. Which of the following is true?\ntag{A}{206}{keph203}{5}{}{Class XI Part 2 \textperiodcentered{} Ch 10}\nsrc{None found in NEET/AIPMT/AIEEE/JEE papers, 2005--2026 (NCERT concept application)}}
{\opt{Buoyancy will be less in water at $0^\circ$C than that in water at $4^\circ$C}{Buoyancy will be more in water at $0^\circ$C than that in water at $4^\circ$C}{Buoyancy in water at $0^\circ$C will be same as that in water at $4^\circ$C}{Buoyancy may be more or less in water at $4^\circ$C depending on the radius of the sphere}}

\QQ{A thermodynamic system undergoes a cyclic process ABC as shown in the diagram. The work done by the system per cycle is
\fig{\begin{tikzpicture}[xscale=0.26,yscale=0.006]
\draw[->] (-1,0) -- (12,0) node[right,lbl]{$V$ (m$^3$)}; \draw[->] (0,-20) -- (0,470) node[above,lbl]{$P$ (N/m$^2$)};
\draw[densely dashed] (0,400) node[left,lbl]{400} -- (5,400) (0,100) node[left,lbl]{100} -- (5,100) (5,100) -- (5,0) node[below,lbl]{5} (10,100) -- (10,0) node[below,lbl]{10};
\draw[thick,postaction={decorate},decoration={markings,mark=at position 0.5 with {\arrow{Stealth}}}] (5,400) -- (5,100);
\draw[thick,postaction={decorate},decoration={markings,mark=at position 0.5 with {\arrow{Stealth}}}] (5,100) -- (10,100);
\draw[thick,postaction={decorate},decoration={markings,mark=at position 0.5 with {\arrow{Stealth}}}] (10,100) -- (5,400);
\node[lbl,above] at (5,400) {A}; \node[lbl,below left] at (5,100) {B}; \node[lbl,right] at (10,100) {C};
\end{tikzpicture}}\ntag{A}{235}{keph204}{10}{}{Class XI Part 2 \textperiodcentered{} Ch 11}\nsrc{NEET 2014 -- thermodynamic system taken round a cyclic process ABC, work done per cycle}}
{\opt{$750\unit{J}$}{$-1250\unit{J}$}{$-750\unit{J}$}{$1250\unit{J}$}}

\QQ{One mole of O$_2$ gas is heated at constant pressure starting at $27^\circ$C. How much energy must be added to the gas as heat to double its volume?\ntag{A}{233}{keph204}{8}{}{Class XI Part 2 \textperiodcentered{} Ch 11}\nsrc{None found in NEET/AIPMT/AIEEE/JEE papers, 2005--2026 (NCERT concept application)}}
{\opt{Zero}{450 R}{750 R}{1050 R}}

\QQ{A piston is performing S.H.M.\ in the vertical direction with a frequency of $0.5\unit{Hz}$. A block of $10\unit{kg}$ is placed on the piston. The maximum amplitude of the system such that the block remains in contact with the piston is\ntag{A}{266}{keph206}{8}{}{Class XI Part 2 \textperiodcentered{} Ch 13}\nsrc{None found in NEET/AIPMT/AIEEE/JEE papers, 2005--2026 (NCERT concept application)}}
{\opt{$1\unit{m}$}{$0.5\unit{m}$}{$1.5\unit{m}$}{$0.1\unit{m}$}}

\QQ{The equation of a stationary wave is $y = 2\sin\left(\dfrac{\pi x}{15}\right)\cos(48\pi t)$. The distance between a node and its next antinode is\ntag{D}{290}{keph207}{13}{}{Class XI Part 2 \textperiodcentered{} Ch 14}}
{\opt{7.5 units}{1.5 units}{22.5 units}{30 units}}

\QQ{In the given circuit, the current through $2\unit{\Omega}$ resistor is
\fig{\begin{circuitikz}[scale=0.55,transform shape,font=\Large]
\draw (0,0) to[battery1,l=$1.2\,\mathrm V$,invert] (0,2.4) to[R,l=$1\,\Omega$] (2.2,2.4) -- (4.4,2.4) to[R,l=$1\,\Omega$] (6.6,2.4) to[R,l=$1\,\Omega$] (6.6,0) to[R,l=$1\,\Omega$] (4.4,0) -- (3.2,0) -- (2.2,0) to[R,l=$2\,\Omega$] (0,0);
\draw (1.6,2.4) to[R,l_=$3\,\Omega$] (1.6,0);
\draw (2.2,2.4) -- (2.2,0); \draw (4.4,2.4) -- (4.4,0);
\draw (2.2,0) to[R,l=$1\,\Omega$] (4.4,2.4); \draw (4.4,0) to[R,l_=$1\,\Omega$] (6.6,2.4);
\end{circuitikz}}\ntag{A}{97}{leph103}{17}{}{Class XII Part 1 \textperiodcentered{} Ch 3}\nsrc{None found in NEET/AIPMT/AIEEE/JEE papers, 2005--2026 (NCERT concept application)}}
{\opt{$0.2\unit{A}$}{$0.3\unit{A}$}{$0.4\unit{A}$}{$0.1\unit{A}$}}

\QQ{Kirchhoff's junction rule is a reflection of\ntag{D}{105}{leph103}{25}{}{Class XII Part 1 \textperiodcentered{} Ch 3}}
{\opt{Conservation of current density vector}{Conservation of energy}{Conservation of momentum}{Conservation of charges}}

\QQ{The variation of terminal potential difference ($V$) with current flowing through a cell is as shown
\fig{\gaxes{$I$ (A)}{$V$ (volt)}{\draw (0,1.5) node[left,lbl]{3 V} -- (1.6,0) node[below,lbl]{6 A};}}
The emf and internal resistance of the cell are\ntag{D}{94}{leph103}{14}{}{Class XII Part 1 \textperiodcentered{} Ch 3}}
{\opt{3\,V, $2\unit{\Omega}$}{3\,V, $0.5\unit{\Omega}$}{6\,V, $2\unit{\Omega}$}{6\,V, $0.5\unit{\Omega}$}}

\QQ{In a potentiometer experiment, the balancing point with a cell is at a length $240\unit{cm}$. On shunting the cell with a resistance of $2\unit{\Omega}$, the balancing length becomes $120\unit{cm}$. The internal resistance of the cell is\ntag{X}{}{}{}{}{potentiometer}}
{\opt{$4\unit{\Omega}$}{$2\unit{\Omega}$}{$1\unit{\Omega}$}{$0.5\unit{\Omega}$}}

\QQ{The magnetic field at the centre `O' in the given figure is
\fig{\begin{tikzpicture}[scale=0.8]
\draw (60:1) arc (60:360:1); \draw[->] (180:1) ++(0,0.05) -- ++(0,-0.15);
\fill (0,0) circle (1.4pt) node[below,lbl]{O};
\draw[densely dashed] (0,0) -- (60:1.25) node[pos=0.55,left,lbl]{$R$}; \draw[densely dashed] (0,0) -- (1.25,0);
\draw (0.35,0) arc (0:60:0.35); \node[lbl] at (0.6,0.32) {$60^\circ$};
\node[lbl] at (-1.25,-0.1) {$I$};
\end{tikzpicture}}\ntag{A}{114}{leph104}{8}{}{Class XII Part 1 \textperiodcentered{} Ch 4}\nsrc{None found in NEET/AIPMT/AIEEE/JEE papers, 2005--2026 (NCERT concept application)}}
{\opt{$\dfrac7{14}\dfrac{\mu_0I}{R}$}{$\dfrac5{12}\dfrac{\mu_0I}{R}$}{$\dfrac3{10}\dfrac{\mu_0I}{R}$}{$\dfrac{\mu_0I}{12R}$}}

\QQ{Two metal plates are separated by $2\unit{cm}$. The potentials of the plates are $-10\unit{V}$ and $+30\unit{V}$. The electric field between the two plates is\ntag{D}{55}{leph102}{11}{}{Class XII Part 1 \textperiodcentered{} Ch 2}}
{\opt{$500\unit{V/m}$}{$1000\unit{V/m}$}{$2000\unit{V/m}$}{$3000\unit{V/m}$}}

\QQ{The equivalent capacitance between A and B is,
\fig{\begin{circuitikz}[scale=0.6,transform shape,font=\Large]
\draw (0,2.6) to[C,l=$50\,\mathrm{pF}$] (4.4,2.6) -- (4.4,1.4) node[circ]{} node[right]{A};
\draw (0,2.6) -- (0,0) to[C,l_=$50\,\mathrm{pF}$] (4.4,0) -- (4.4,0.6) node[circ]{} node[right]{B};
\draw (0,1.4) to[C,l_=$100\,\mathrm{pF}$] (2.2,1.4) to[C,l_=$100\,\mathrm{pF}$] (4.4,1.4);
\end{circuitikz}}\ntag{A}{72}{leph102}{28}{}{Class XII Part 1 \textperiodcentered{} Ch 2}\nsrc{None found in NEET/AIPMT/AIEEE/JEE papers, 2005--2026 (NCERT concept application)}}
{\opt{$50\unit{pF}$}{$\dfrac{100}3\unit{pF}$}{$150\unit{pF}$}{$300\unit{pF}$}}

\QQ{A capacitor of capacitance C charged by an amount Q is connected in parallel with an uncharged capacitor of capacitance 2C. The final charges on the capacitors are\ntag{A}{73}{leph102}{29}{}{Class XII Part 1 \textperiodcentered{} Ch 2}\nsrc{None found in NEET/AIPMT/AIEEE/JEE papers, 2005--2026 (NCERT concept application)}}
{\opt{$\dfrac Q2,\ \dfrac Q2$}{$\dfrac Q4,\ \dfrac{3Q}4$}{$\dfrac Q3,\ \dfrac{2Q}3$}{$\dfrac Q5,\ \dfrac{4Q}5$}}

\QQ{Though the electron drift velocity is small and electron charge is very small, a conductor can carry an appreciably large current because\ntag{D}{88}{leph103}{8}{ \textperiodcentered{} Example 3.2}{Class XII Part 1 \textperiodcentered{} Ch 3}}
{\opt{electron number density is very large}{drift velocity of electron is very large}{electron number density depends on temperature}{relaxation time is small}}

\QQ{Masses of three wires of copper are in the ratio 1:3:5 and their lengths are in the ratio 5:3:1. The ratio of their electrical resistance are\ntag{A}{84}{leph103}{4}{}{Class XII Part 1 \textperiodcentered{} Ch 3}\nsrc{None found in NEET/AIPMT/AIEEE/JEE papers, 2005--2026 (NCERT concept application)}}
{\opt{1:3:5}{5:3:1}{1:15:125}{125:15:1}}

\QQ{A certain charge 2Q is divided at first into two parts $q_1$ and $q_2$. Later the charges are placed at a certain distance. If the force of interaction between two charges is maximum then $\dfrac Q{q_1} =$ \rule{1cm}{0.4pt}.\ntag{A}{8}{leph101}{8}{}{Class XII Part 1 \textperiodcentered{} Ch 1}\nsrc{None found in NEET/AIPMT/AIEEE/JEE papers, 2005--2026 (NCERT concept application)}}
{\opt{4}{2}{1}{0.5}}

\QQ{A particle of mass $m$ and charge $q$ is placed at rest in uniform electric field $E$ and then released. The kinetic energy attained by the particle after moving a distance $y$ is\ntag{D}{59}{leph102}{15}{}{Class XII Part 1 \textperiodcentered{} Ch 2}}
{\opt{$qEy^2$}{$qE^2y$}{$qEy$}{$q^2Ey$}}

\QQ{An electric dipole is kept in non-uniform electric field. It generally experiences\ntag{D}{28}{leph101}{28}{}{Class XII Part 1 \textperiodcentered{} Ch 1}}
{\opt{A force and torque}{A force but not a torque}{A torque but not a force}{Neither a force nor a torque}}

\QQ{The figure gives the electric potential $V$ as a function of distance through four regions on $x$-axis. Which of the following is true for the magnitude of the electric field E in these regions?
\fig{\begin{tikzpicture}[xscale=0.6,yscale=0.5]
\draw[->] (-0.3,0) -- (5.8,0) node[right,lbl]{$x$ (m)}; \draw[->] (0,-0.3) -- (0,4.8) node[above,lbl]{$V$ (volt)};
\foreach \x in {1,...,5} \node[lbl,below] at (\x,0) {\x}; \node[lbl,left] at (0,2) {2}; \node[lbl,left] at (0,4) {4};
\draw[thick] (0,2) -- (1,2) -- (2,4) -- (4,4) -- (4.6,0);
\foreach \x in {1,2,4} \draw[densely dashed] (\x,0) -- (\x,4);
\node[font=\scriptsize] at (0.5,1.0) {\textcircled{\tiny A}}; \node[font=\scriptsize] at (1.5,1.0) {\textcircled{\tiny B}}; \node[font=\scriptsize] at (3,1.0) {\textcircled{\tiny C}}; \node[font=\scriptsize] at (4.75,0.6) {\textcircled{\tiny D}};
\end{tikzpicture}}\ntag{D}{55}{leph102}{11}{}{Class XII Part 1 \textperiodcentered{} Ch 2}}
{\opt{$E_A>E_B>E_C>E_D$}{$E_A = E_C$ and $E_B<E_D$}{$E_B = E_D$ and $E_A<E_C$}{$E_A<E_B<E_C<E_D$}}

\QQ{A system of two charges separated by a certain distance apart stores electrical potential energy. If the distance between them is increased, the potential energy of the system,\ntag{D}{56}{leph102}{12}{}{Class XII Part 1 \textperiodcentered{} Ch 2}}
{\opt{increases in any case}{decreases in any case}{may increase or decrease}{remains the same}}

%%BEGIN-ANALYSIS
\end{multicols}
\par\needspace{16\baselineskip}\vspace{14pt}
\begin{center}{\pgfornament[width=3.2cm,color=gilt]{87}}\par
{\large\color{navy}\scshape KCET 2019 \textperiodcentered{} Answer Key}\end{center}
{\small\setlength{\tabcolsep}{1pt}\renewcommand{\arraystretch}{1.5}
\begin{tabularx}{\linewidth}{@{}*{10}{>{\centering\arraybackslash}X}@{}}\toprule
\textcolor{gray}{\textbf{1}\,(1)$^{\dagger}$} & \textbf{2}\,(2) & \textbf{3}\,(4) & \textbf{4}\,(3) & \textcolor{gray}{\textbf{5}\,(3)$^{\dagger}$} & \textcolor{gray}{\textbf{6}\,(1)$^{\dagger}$} & \textbf{7}\,(4) & \textcolor{gray}{\textbf{8}\,(3)$^{\dagger}$} & \textcolor{gray}{\textbf{9}\,(3)$^{\dagger}$} & \textcolor{gray}{\textbf{10}\,(1)$^{\dagger}$}\\
\textbf{11}\,(1) & \textbf{12}\,(3) & \textbf{13}\,(2) & \textbf{14}\,(2) & \textbf{15}\,(3) & \textbf{16}\,(1) & \textbf{17}\,(4) & \textbf{18}\,(1) & \textcolor{gray}{\textbf{19}\,(2)$^{\dagger}$} & \textbf{20}\,(2)\\
\textcolor{gray}{\textbf{21}\,(3)$^{\dagger}$} & \textbf{22}\,(2) & \textcolor{gray}{\textbf{23}\,(3)$^{\dagger}$} & \textbf{24}\,(1) & \textbf{25}\,(2)$^{*}$ & \textcolor{gray}{\textbf{26}\,(1)$^{*}$$^{\dagger}$} & \textcolor{gray}{\textbf{27}\,(1)$^{\dagger}$} & \textbf{28}\,(2) & \textbf{29}\,(2) & \textbf{30}\,(4)\\
\textbf{31}\,(2) & \textbf{32}\,(3) & \textbf{33}\,(1) & \textbf{34}\,(1) & \textbf{35}\,(1) & \textbf{36}\,(2) & \textbf{37}\,(4) & \textbf{38}\,(2) & \textbf{39}\,(4) & \textbf{40}\,(3)\\
\textbf{41}\,(1) & \textbf{42}\,(3) & \textbf{43}\,(4) & \textbf{44}\,(1) & \textbf{45}\,(1) & \textbf{46}\,(3) & \textbf{47}\,(4) & \textbf{48}\,(2) & \textcolor{gray}{\textbf{49}\,(2)$^{\dagger}$} & \textbf{50}\,(2)\\
\textbf{51}\,(3) & \textbf{52}\,(2) & \textbf{53}\,(3) & \textbf{54}\,(1) & \textbf{55}\,(4) & \textbf{56}\,(3) & \textbf{57}\,(3) & \textbf{58}\,(1) & \textbf{59}\,(2) & \textbf{60}\,(3)\\\bottomrule\end{tabularx}}\par\smallskip
{\footnotesize Answers solved by the editor; (1)--(4) are the printed options. {\color{gray}Grey$^{\dagger}$}: deleted portion (answer given for completeness). $^{*}$: see note.\par
\par\smallskip\textbf{Q25.} NCERT: saturation current is proportional to intensity. Strictly, doubling both frequency and intensity leaves the number of photons per second unchanged.
\par\smallskip\textbf{Q26.} The data give $B = 4\pi^2 I/(MT^2) \approx 1.6\times10^{-3}$\,T; option (1) is the intended one.
}\par
\clearpage
\begin{center}{\pgfornament[width=5cm,color=gilt]{87}}\par
{\fontsize{20}{24}\selectfont\color{navy}\scshape KCET 2019 \textperiodcentered{} Paper Analysis}\par
{\color{gilt}\pgfornament[width=5cm,color=gilt,symmetry=h]{87}}\end{center}
\vspace{4pt}
\begin{multicols}{2}
\textbf{\color{crimson}1. Direct vs Applied}\par\smallskip
{\footnotesize Checked question by question against the four 2026-27 NCERT books (XI Part 1 \& 2, XII Part 1 \& 2). Out of 60, 48 questions are in the present syllabus; percentages are of these.}\par\smallskip
\begin{tabular}{@{}lrr@{}}\toprule
Direct from NCERT & \textbf{20} & 42\%\\
\quad text, summary, table or figure & 18 &\\
\quad NCERT solved Example & 1 &\\
\quad NCERT Exercise & 1 &\\
Applied (concept in NCERT) & \textbf{28} & 58\%\\\midrule
Analysed & 48 &\\
Deleted portion (not analysed) & 12 &\\\midrule
\textbf{Total} & \textbf{60} &\\\bottomrule\end{tabular}\par\smallskip
{\footnotesize Direct: Q 2, 4, 7, 11, 15, 16, 20, 22, 25, 31, 37, 45, 47, 48, 51, 54, 57, 58, 59, 60.\\ From NCERT Examples: Q 54 (Example 3.2).\\ From NCERT Exercises: Q 22 (Exercise 10.5).}\par\medskip
\textbf{\color{crimson}2. Where the applied questions resemble earlier exams}\par\smallskip
Of the 28 applied questions, \textbf{1} closely resemble a question from an earlier national exam:
NEET/AIPMT/AIIMS 1, JEE Main/Advanced 0, AIEEE 0. The other 27 (Q 3, 12, 13, 14, 17, 18, 24, 28, 29, 30, 32, 33, 34, 35, 36, 38, 39, 40, 41, 43, 44, 46, 50, 52, 53, 55, 56) had no close match in the NEET, AIPMT, AIIMS, AIEEE and JEE papers of 2005--2026 that were searched; they are straightforward applications of NCERT concepts.\par
{\footnotesize Strongest repeats: Q42 (NEET 2014).}\par\medskip
\textbf{\color{crimson}3. Difficulty level}\par\smallskip
\begin{tabular}{@{}lr@{}}\toprule
Easy & 29\\ Moderate & 18\\ Hard & 1\\\bottomrule\end{tabular}\par\smallskip
{\footnotesize Hard: Q 39. Overall: \textbf{easy to moderate} and heavily conceptual; many questions are NCERT Exemplar MCQs. Q39 (image speed for a moving object in front of a concave mirror) is the hardest. Class XI: 15 questions, Class XII: 33 (analysed questions only). (Difficulty is the editor's assessment.)}\par\medskip
\columnbreak
\textbf{\color{crimson}4. Chapter-wise distribution}\par\smallskip
{\footnotesize\begin{tabular}{@{}p{4.7cm}c@{\,}r@{\,}r@{\,}r@{}}\toprule
Chapter & Cl. & Qs & Dir & App\\\midrule
Units \& Measurement & XI & 1 & 1 & 0\\
Motion in a Straight Line & XI & 1 & 0 & 1\\
Motion in a Plane & XI & 1 & 0 & 1\\
Laws of Motion & XI & 1 & 0 & 1\\
Work, Energy \& Power & XI & 1 & 1 & 0\\
Rotational Motion & XI & 2 & 1 & 1\\
Gravitation & XI & 1 & 0 & 1\\
Mech.\ Prop.\ of Solids & XI & 1 & 0 & 1\\
Mech.\ Prop.\ of Fluids & XI & 1 & 0 & 1\\
Thermal Prop.\ of Matter & XI & 1 & 0 & 1\\
Thermodynamics & XI & 2 & 0 & 2\\
Oscillations & XI & 1 & 0 & 1\\
Waves & XI & 1 & 1 & 0\\
Electric Charges \& Fields & XII & 2 & 1 & 1\\
Electrostatic Potential \& Capacitance & XII & 6 & 4 & 2\\
Current Electricity & XII & 6 & 4 & 2\\
Moving Charges \& Magnetism & XII & 3 & 1 & 2\\
Magnetism \& Matter & XII & 1 & 1 & 0\\
Electromagnetic Induction & XII & 1 & 0 & 1\\
Alternating Current & XII & 3 & 0 & 3\\
Electromagnetic Waves & XII & 1 & 1 & 0\\
Ray Optics & XII & 3 & 0 & 3\\
Wave Optics & XII & 1 & 1 & 0\\
Dual Nature & XII & 2 & 1 & 1\\
Atoms & XII & 3 & 1 & 2\\
Nuclei & XII & 1 & 1 & 0\\
\bottomrule\end{tabular}}\par\medskip
\textbf{\color{crimson}5. Teaching takeaways}\par\smallskip
{\footnotesize\begin{itemize}[leftmargin=1em,nosep]
\item Exercise 10.5 (intensity $K$ at path difference $\lambda$, find it at $\lambda/3$) was asked word for word; Example 3.2 (large electron number density) appeared as a statement question.
\item NCERT text asked: doubling the turns of a galvanometer (p.\,131), semiconductor conductivity rising with temperature (p.\,91), force and torque on a dipole in a non-uniform field (p.\,28), sign of the potential energy of two charges (p.\,56).
\item Twelve questions are on deleted portions (cyclotron, coercivity, mean life, transistor, logic gates, antenna, decay, Brewster's law, Doppler effect, needle oscillation, toroid, potentiometer); KCET 2019 was set on the pre-2023 syllabus.
\end{itemize}}\par\medskip
\textbf{\color{crimson}6. Deleted portions}\par\smallskip
\textbf{12} of the 60 questions are on topics removed from the 2026-27 NCERT (rationalised syllabus). They are marked \textsc{Deleted Portion} and are not analysed: cyclotron (Q 1); coercivity (hysteresis) (Q 5); radioactive decay law (mean life) (Q 6); transistor amplifier (Q 8); digital electronics (logic gates) (Q 9); antenna size (communication systems) (Q 10); alpha- and beta-decay (Q 19); Brewster's law (Q 21); Doppler effect for light (Q 23); oscillation period of a magnetic needle (Q 26); toroid (Q 27); potentiometer (Q 49).\par
\end{multicols}
%%END-ANALYSIS







